\documentclass[12pt,a4paper]{article}
\usepackage{ctex}
\usepackage{geometry}
\usepackage{graphicx}
\usepackage{xcolor}
\usepackage{float}
\usepackage{hyperref}
\usepackage{fancyhdr}

\geometry{left=2cm,right=2cm,top=2.5cm,bottom=2.5cm}
\setlength{\headheight}{15pt}

\definecolor{accent}{RGB}{59,130,246}

\hypersetup{
  colorlinks=true,
  linkcolor=accent,
  urlcolor=accent,
  pdftitle={Dionysus 前端组件与界面图册},
  pdfauthor={FuyuuKu樱}
}

\pagestyle{fancy}
\fancyhf{}
\fancyhead[L]{Dionysus 前端组件与界面图册}
\fancyhead[R]{\thepage}
\renewcommand{\headrulewidth}{0.4pt}

\title{\textbf{Dionysus 前端组件与界面图册}\\[0.3em]\large 主要界面 \& React 组件状态参考}
\author{FuyuuKu樱}
\date{版本 v0.2.0 \quad 2026-06-24}

\begin{document}

\maketitle
\thispagestyle{empty}

\vfill
\begin{center}
\textit{让 AI Agent 拥有一个会动、会说的看板娘}
\end{center}
\vfill

\newpage
\tableofcontents
\newpage

% ============================================================
\section{布局与导航}
% ============================================================

\begin{figure}[H]
  \centering
  \includegraphics[width=\textwidth]{../qa_screenshots/component_gallery/01_layout_default.png}
  \caption{默认三栏布局：左侧会话与导航、中间聊天区、右侧角色与工具面板。}
\end{figure}

\begin{figure}[H]
  \centering
  \includegraphics[width=\textwidth]{../qa_screenshots/component_gallery/02_palette.png}
  \caption{角色与调色盘切换：通过顶部标签快速切换不同角色主题。}
\end{figure}

\begin{figure}[H]
  \centering
  \includegraphics[width=\textwidth]{../qa_screenshots/component_gallery/06_session_menu.png}
  \caption{会话管理菜单：对当前会话进行重命名、清空、删除等操作。}
\end{figure}

% ============================================================
\section{主要页面}
% ============================================================

\begin{figure}[H]
  \centering
  \includegraphics[width=\textwidth]{../qa_screenshots/component_gallery/03_persona.png}
  \caption{角色设置页：管理角色 YAML 配置、Live2D 模型与语料。}
\end{figure}

\begin{figure}[H]
  \centering
  \includegraphics[width=\textwidth]{../qa_screenshots/component_gallery/04_settings.png}
  \caption{系统设置页：配置连接地址、主题、快捷键与行为选项。}
\end{figure}

\begin{figure}[H]
  \centering
  \includegraphics[width=\textwidth]{../qa_screenshots/component_gallery/05_qr_popup.png}
  \caption{二维码弹窗：展示移动端配对或 Web 端访问入口。}
\end{figure}

% ============================================================
\section{聊天消息}
% ============================================================

\begin{figure}[H]
  \centering
  \includegraphics[width=\textwidth]{../qa_screenshots/component_gallery/07_chat_messages.png}
  \caption{消息气泡：用户消息与 Agent 回复的分层展示。}
\end{figure}

\begin{figure}[H]
  \centering
  \includegraphics[width=\textwidth]{../qa_screenshots/component_gallery/08_chat_thinking.png}
  \caption{思考中状态：Agent 在生成正式回复前展示思考过程。}
\end{figure}

\begin{figure}[H]
  \centering
  \includegraphics[width=\textwidth]{../qa_screenshots/component_gallery/09_chat_toolcall.png}
  \caption{工具调用：右侧工具面板与全局 HUD 同时显示运行中的工具。}
\end{figure}

% ============================================================
\section{交互组件}
% ============================================================

\begin{figure}[H]
  \centering
  \includegraphics[width=\textwidth]{../qa_screenshots/component_gallery/10_options_button_group.png}
  \caption{选项按钮组：用户可在聊天底部一键选择预设方案。}
\end{figure}

\begin{figure}[H]
  \centering
  \includegraphics[width=\textwidth]{../qa_screenshots/component_gallery/11_options_dropdown.png}
  \caption{选项下拉框：以紧凑下拉形式呈现多个可选项。}
\end{figure}

\begin{figure}[H]
  \centering
  \includegraphics[width=\textwidth]{../qa_screenshots/component_gallery/12_options_card_list.png}
  \caption{选项卡片列表：以卡片形式展示带描述的选项。}
\end{figure}

\begin{figure}[H]
  \centering
  \includegraphics[width=\textwidth]{../qa_screenshots/component_gallery/13_todos.png}
  \caption{待办面板：展示任务清单与完成进度。}
\end{figure}

\begin{figure}[H]
  \centering
  \includegraphics[width=\textwidth]{../qa_screenshots/component_gallery/14_streaming_status.png}
  \caption{流式状态：显示 Agent 当前的执行阶段与进度条。}
\end{figure}

% ============================================================
\section{角色陪伴}
% ============================================================

\begin{figure}[H]
  \centering
  \includegraphics[width=\textwidth]{../qa_screenshots/component_gallery/15_companion_dialog.png}
  \caption{角色对话框：看板娘以角色口吻主动播报状态与问候。}
\end{figure}

\end{document}
